\section*{Summary}
\addcontentsline{toc}{section}{Summary}

\begin{itemize}
  \item \code{Interface} makes the application testable; \code{ByteTransport} makes \textbf{the
    driver} testable. Without the second seam, \code{Spi} runs for the first time on hardware
    somebody else is debugging at the same time.
  \item The transport is deliberately dumb: \code{select()}, \code{deselect()},
    \code{transfer()}. It does not know that a transaction is five bytes long --- the framing
    belongs to the protocol.
  \item The register constants are \textbf{indices}, not offsets, and the masks are written as
    shifts so that they can be compared line by line against the register map.
  \item \code{readReg()} discards the byte clocked in during the command byte. Four out of five
    contribute.
  \item \textbf{Cast before the shift.} Otherwise a \code{std::uint8_t} is promoted to \code{int},
    and a byte with bit 7 set shifted 24 steps writes into the sign bit.
  \item The helper functions are private and contain no logic. All the meaning lives in the layer
    above.
\end{itemize}
